% Schémas de Lewis des molécules et des ions du programme (L'essentiel, chapitre 6)
{\renewcommand{\arraystretch}{1}\setlength{\tabcolsep}{2pt}%
\newcommand{\nomL}[1]{{\otLight\fontsize{7.5}{9}\selectfont\color{otMuted}#1}}%
\begin{tabular}{*{5}{>{\centering\arraybackslash}m{30mm}}}
\chemfig{H-H} & \chemfig{\charge{135=\|,225=\|}{O}=\charge{45=\|,315=\|}{O}} & \chemfig{\charge{180=\|}{N}~\charge{0=\|}{N}} &
\chemfig{H-\charge{0=\|,90=\|,270=\|}{Cl}} & \chemfig{\charge{135=\|,225=\|}{O}=C=\charge{45=\|,315=\|}{O}} \\[1mm]
\nomL{dihydrogène} & \nomL{dioxygène} & \nomL{diazote} & \nomL{chlorure d'hydrogène} & \nomL{dioxyde de carbone} \\[4mm]
\chemfig{H-\charge{90=\|,270=\|}{O}-H} & \chemfig{H-\charge{90=\|}{N}(-[6]H)-H} & \chemfig{H-C(-[2]H)(-[6]H)-H} &
\chemfig{H-\charge{90=\|,45:3pt=\cplus}{O}(-[6]H)-H} & \chemfig{H-\charge{45:3pt=\cplus}{N}(-[2]H)(-[6]H)-H} \\[1mm]
\nomL{eau} & \nomL{ammoniac} & \nomL{méthane} & \nomL{ion oxonium} & \nomL{ion ammonium} \\[4mm]
\chemfig{\charge{0=\"}{H}}\,\raisebox{1.2ex}{\cplus} & \chemfig{Na}\,\raisebox{1.2ex}{\cplus} &
\chemfig{\charge{0=\|,90=\|,180=\|,270=\|,45:3pt=\cmoins}{Cl}} &
\chemfig{\charge{0=\|,90=\|,180=\|,270=\|,45:3pt=\cmoins,135:3pt=\cmoins}{O}} &
\chemfig{\charge{90=\|,180=\|,270=\|,135:3pt=\cmoins}{O}-H} \\[1mm]
\nomL{ion hydrogène (lacune)} & \nomL{ion sodium} & \nomL{ion chlorure} & \nomL{ion oxyde} & \nomL{ion hydroxyde} \\
\end{tabular}}
\par\smallskip\legende{Un trait entre deux atomes est un doublet liant ; un trait autour d'un atome, un doublet non liant ; un rectangle vide, une lacune électronique.}
